% BridgeIR — paper draft.
% Convention: paper prose lives in docs/ (this file + sections/).
% Working notes, decisions, scratch material live in docs/scratch/.
\documentclass[11pt]{article}

\usepackage[T1]{fontenc}
\usepackage[margin=1in]{geometry}
\usepackage{amsmath,amssymb}
\usepackage{booktabs}
\usepackage{tabularx}
\usepackage{array}
\usepackage{ragged2e}
\usepackage{graphicx}
\usepackage{xcolor}
\usepackage[numbers,sort&compress]{natbib}
\usepackage{hyperref}

\newcommand{\mlir}{MLIR}

\title{BridgeIR: An Agent-Operated, Runtime-Extensible Multi-Backend IR Stack}
\author{Hocky Yudhiono \\ \small (draft --- authors TBD)}
\date{\today}

\begin{document}
\maketitle

\begin{abstract}
Draft. BridgeIR is a multi-backend intermediate-representation stack whose extension
surface --- pipelines, pass definitions, verifiers, and invariants --- is expressed as
data that an automated agent can compose and load at runtime. We adopt \mlir{} as the
IR substrate and contribute the layers it lacks: a tiered extension model with precise
rebuild boundaries, a runtime control plane, and first-class invariants that make
agent-authored transformations trustworthy.
{\color{red}\textbf{TODO:}} sharpen thesis, state contributions, results.
\end{abstract}

\section{Introduction}
\label{sec:intro}

Compiler pipelines are increasingly built and operated by automated agents, yet the
frameworks they must extend were designed for humans editing C++ and waiting for
rebuilds. Adding a pass, an operation definition, a verifier, or an invariant to an
\mlir{}-based stack today means editing TableGen and C++ sources, rebuilding the tool,
and restarting the process. That authoring loop --- not the IR data structure --- is
the bottleneck.

Our thesis is that the IR \emph{core} is not the right place to compete: \mlir{} has
won the core, and systems that try to replace it lose the ecosystem while chasing a
moving target. What is missing around the core are (i) declarative extension tiers
that an agent can author as data, (ii) a control plane that discovers, validates, and
loads those extensions at runtime, and (iii) first-class invariants that make
agent-authored transformations trustworthy.

\textbf{Contributions (draft).}
\begin{itemize}
  \item A tiered extension model for agent-authored compiler extensions --- from
        pipeline manifests and interpreted transformations up to runtime operation
        definitions and compiled plugins --- with precise rebuild boundaries.
        {\color{red}\textbf{TODO}}
  \item A runtime control plane: a manifest registry, plugin loader, and structured
        (machine-readable) diagnostics.
        {\color{red}\textbf{TODO}}
  \item Invariants as data: per-stage contracts checked at runtime.
        {\color{red}\textbf{TODO:} decide SMT-checked vs.\ translation-validation depth.}
  \item An evaluation protocol for \emph{agent operability}: task success, time and
        cost to first successful extension, error recovery, versus a vanilla
        out-of-tree \mlir{} workflow.
        {\color{red}\textbf{TODO:} define tasks, baselines, metrics.}
\end{itemize}

\section{Positioning: Where \mlir{} Ends}
\label{sec:positioning}

We frame the design space as a set of axes for \mlir{}~\cite{lattner2021mlir}. For
each axis we ask what is available today and where systems that go \emph{beyond} it
live. ``Bigger'' here does not mean a single system that dominates \mlir{} --- no such
system exists today --- but rather the research or engineering direction that pushes
that particular axis past what \mlir{} offers. Table~\ref{tab:landscape} states the
comparison.

\begin{table}[t]
  \centering
  \small
  \begin{tabularx}{\textwidth}{@{}>{\raggedright\arraybackslash}p{2.6cm} >{\raggedright\arraybackslash}X >{\raggedright\arraybackslash}X@{}}
  \toprule
  \textbf{Axis} & \textbf{\mlir{} today} & \textbf{Where ``bigger'' lives} \\
  \midrule
  Multi-level structure &
  Top form: a metalanguage; multi-level IRs are cheap &
  --- (no rival supersedes it) \\
  Declarative definitions (ops/types/attrs) &
  Top form: ODS/TableGen generates builders, accessors, verifiers &
  --- \\
  Interpreted transformations &
  Young: PDL patterns and transform-dialect scripts &
  E-graph rewriting: egg~\cite{willsey2021egg},
  egglog~\cite{zhang2023better}; tensor superoptimization:
  TASO~\cite{jia2019taso} \\
  Rewrite search and scheduling &
  Greedy, order-sensitive; fixpoints via ad-hoc pass orders &
  Equality saturation; relational e-matching~\cite{zhang2022relational};
  rule inference~\cite{nandi2021ruler} \\
  Per-compilation specialization &
  Static pipelines; passes fixed at build time; schedules left to backends &
  Autotuned schedules (TVM~\cite{chen2018tvm}, Triton~\cite{tillet2019triton});
  peephole synthesis (STOKE~\cite{schkufza2013stoke}); generative compilation
  (SPIRAL~\cite{puschel2005spiral}); staged specialization
  (Futamura~\cite{futamura1971partial}) \\
  Semantics and guarantees &
  Local per-op C++ verifiers; no formal semantics &
  Translation validation (Alive2)~\cite{lopes2021alive2}; mechanized semantics
  (Vellvm)~\cite{zhao2012vellvm}; SMT-backed contracts \\
  Incrementality &
  Whole-pass recomputation &
  Demand-driven compilation: Adapton~\cite{hammer2014adapton},
  Salsa~\cite{salsa} \\
  Runtime extensibility &
  Second-class: \texttt{ExtensibleDialect} with C++ verify callbacks;
  \texttt{.so} plugins with exact-ABI requirements &
  Data-first extension registries; sandboxed plugin hosts (WASM) \\
  Agent operability &
  Weak: extension means editing \texttt{.td}/C++ and rebuilding &
  Plan $\rightarrow$ build $\rightarrow$ load loops; machine-readable
  pass/dialect descriptions; structured diagnostics \\
  Stability and interchange &
  No core stability; dialect-level fixes (StableHLO)~\cite{stablehlo} &
  Versioned dialects with serialization discipline \\
  Schedule and hardware semantics &
  Partial: affine, linalg, CIRCT &
  Schedule-first IRs: Halide~\cite{ragankelley2013halide},
  TVM TensorIR~\cite{zheng2023tensorir}, ISL schedule trees \\
  \bottomrule
  \end{tabularx}
  \caption{Axis-by-axis comparison. ``Bigger'' names the direction that pushes an axis
  past \mlir{}, not a single competing system.}
  \label{tab:landscape}
\end{table}

Reading the table: the first two rows are the foundation we inherit unchanged. The
middle rows constitute the design space BridgeIR occupies --- interpreted
transformations, trustworthy rewriting, per-compilation specialization, invariants,
incrementality, runtime extensibility, and agent operability. The last rows we adopt
as discipline (versioned dialects, stable serialization) and per-backend work
(schedule and hardware semantics).
{\color{red}\textbf{TODO:} add a per-axis ``BridgeIR stance'' once the design section
stabilizes.}


% Planned sections (write next):
% \input{sections/03-background}
% \input{sections/04-design}
% \input{sections/05-implementation}
% \input{sections/06-evaluation}
% \section{Related Work}
\label{sec:related}

\paragraph{Compiler IRs.} LLVM~\cite{lattner2004llvm} and GCC's
GIMPLE~\cite{merrill2003gimple} established CFG+SSA as the production mid-end form;
Swift's SIL~\cite{swiftsil} adds ownership and staged canonicalization; Rust's
MIR~\cite{mir} is designed for borrow-checking and interpreted by Miri, demonstrating
IR-as-executable. Functional compilers contribute typed cores and the many-small-IRs
discipline: GHC Core/STG~\cite{peytonjones1992stg}, Core Erlang~\cite{coreerlang}, and
the nanopass framework~\cite{keep2013nanopass}, which auto-generates a verifier for
every tiny stage language. BridgeIR borrows the small-stage philosophy but makes stage
definitions and their checks data, and keeps node identity stable across stages.

\paragraph{ML compilers and tensor IRs.} XLA's HLO and StableHLO~\cite{stablehlo}
treat stability as an explicit feature; TVM~\cite{chen2018tvm} and
TensorIR~\cite{zheng2023tensorir} make schedules first-class programs; Triton
encodes hardware layouts in IR types~\cite{tillet2019triton}; Halide separates
algorithm from schedule~\cite{ragankelley2013halide}; PyTorch's FX~\cite{pytorchfx}
keeps the graph as data structures in the host language. TASO~\cite{jia2019taso},
egg~\cite{willsey2021egg}, egglog~\cite{zhang2023better}, Ruler~\cite{nandi2021ruler},
and relational e-matching~\cite{zhang2022relational} contribute equivalence-based
rewriting and rule inference. BridgeIR adopts schedules-as-data and
equality-friendly rewriting, but targets whole-pipeline agent operation rather than
single-kernel optimization.

\paragraph{Verification and provenance.} Verification-oriented IRs — Boogie and its
descendants~\cite{barnett2005boogie} — place contracts at the center of the
representation. Production IR verifiers (JVM~\cite{jvmspec}, wasm~\cite{haas2017wasm},
eBPF~\cite{ebpf}, SPIR-V~\cite{spirvspec}) demonstrate strong static checking
precedents, and database optimizers show how to record derivation
history~\cite{graefe1995cascades}. BridgeIR lifts both ideas to every transformation:
transactional commits with invariant gates, and deltas with provenance as first-class
values.

\paragraph{Agent-driven compilation.} Recent work pairs LLM agents with compilers:
AlphaEvolve~\cite{alphaevolve} evolves kernels and algorithms with verification
loops; KernelBench~\cite{kernelbench} benchmarks LLM kernel generation; Meta's LLM
Compiler~\cite{llmcompiler} fine-tunes models for LLVM IR tasks;
ComPilot~\cite{compilot} assists compiler debugging; Compiler-R1~\cite{compilerr1}
performs reinforcement-learning-driven optimization tuning. All of these operate on
IRs designed for human pass authors. To our knowledge, no system redesigns the IR
itself around agent affordances.

\paragraph{Positioning.} BridgeIR combines four elements we did not find together in
any surveyed system: stable node identity that survives passes and serialization;
definitions as runtime data; transactional deltas with provenance; and a canonical,
token-lean text format. Each element has precedents — exchange-format naming, ONNX
schemas, TVM runtime op registration, Cascades-style derivation tracking, WAT/PTX
text canons — but their conjunction, targeted at agents, is new.

% \input{sections/08-discussion}
% \input{sections/09-conclusion}

\bibliographystyle{plainnat}
\bibliography{refs}

\end{document}
